% Options for packages loaded elsewhere
% Options for packages loaded elsewhere
\PassOptionsToPackage{unicode}{hyperref}
\PassOptionsToPackage{hyphens}{url}
%
\documentclass[
  ignorenonframetext,
  aspectratio=169,
  russian,
  aspectratio=169]{beamer}
\newif\ifbibliography
\usepackage{pgfpages}
\setbeamertemplate{caption}[numbered]
\setbeamertemplate{caption label separator}{: }
\setbeamercolor{caption name}{fg=normal text.fg}
\beamertemplatenavigationsymbolshorizontal
% remove section numbering
\setbeamertemplate{part page}{
  \centering
  \begin{beamercolorbox}[sep=16pt,center]{part title}
    \usebeamerfont{part title}\insertpart\par
  \end{beamercolorbox}
}
\setbeamertemplate{section page}{
  \centering
  \begin{beamercolorbox}[sep=12pt,center]{section title}
    \usebeamerfont{section title}\insertsection\par
  \end{beamercolorbox}
}
\setbeamertemplate{subsection page}{
  \centering
  \begin{beamercolorbox}[sep=8pt,center]{subsection title}
    \usebeamerfont{subsection title}\insertsubsection\par
  \end{beamercolorbox}
}
% Prevent slide breaks in the middle of a paragraph
\widowpenalties 1 10000
\raggedbottom
\AtBeginPart{
  \frame{\partpage}
}
\AtBeginSection{
  \ifbibliography
  \else
    \frame{\sectionpage}
  \fi
}
\AtBeginSubsection{
  \frame{\subsectionpage}
}
\usepackage{iftex}
\ifPDFTeX
  \usepackage[T1]{fontenc}
  \usepackage[utf8]{inputenc}
  \usepackage{textcomp} % provide euro and other symbols
\else % if luatex or xetex
  \usepackage{unicode-math} % this also loads fontspec
  \defaultfontfeatures{Scale=MatchLowercase}
  \defaultfontfeatures[\rmfamily]{Ligatures=TeX,Scale=1}
\fi
\usepackage{lmodern}

\usefonttheme{serif} % use mainfont rather than sansfont for slide text
\ifPDFTeX\else
  % xetex/luatex font selection
  \setmainfont[]{DejaVu Serif}
  \setsansfont[]{DejaVu Sans}
  \setmonofont[]{DejaVu Sans Mono}
\fi
% Use upquote if available, for straight quotes in verbatim environments
\IfFileExists{upquote.sty}{\usepackage{upquote}}{}
\IfFileExists{microtype.sty}{% use microtype if available
  \usepackage[]{microtype}
  \UseMicrotypeSet[protrusion]{basicmath} % disable protrusion for tt fonts
}{}


\usepackage{longtable,booktabs,array}
\usepackage{calc} % for calculating minipage widths
\usepackage{caption}
% Make caption package work with longtable
\makeatletter
\def\fnum@table{\tablename~\thetable}
\makeatother
\usepackage{graphicx}
\makeatletter
\newsavebox\pandoc@box
\newcommand*\pandocbounded[1]{% scales image to fit in text height/width
  \sbox\pandoc@box{#1}%
  \Gscale@div\@tempa{\textheight}{\dimexpr\ht\pandoc@box+\dp\pandoc@box\relax}%
  \Gscale@div\@tempb{\linewidth}{\wd\pandoc@box}%
  \ifdim\@tempb\p@<\@tempa\p@\let\@tempa\@tempb\fi% select the smaller of both
  \ifdim\@tempa\p@<\p@\scalebox{\@tempa}{\usebox\pandoc@box}%
  \else\usebox{\pandoc@box}%
  \fi%
}
% Set default figure placement to htbp
\def\fps@figure{htbp}
\makeatother



\ifLuaTeX
\usepackage[bidi=basic,provide=*]{babel}
\else
\usepackage[bidi=default,provide=*]{babel}
\fi
\ifPDFTeX
\else
\babelfont{rm}[]{DejaVu Serif}
\fi
% get rid of language-specific shorthands (see #6817):
\let\LanguageShortHands\languageshorthands
\def\languageshorthands#1{}


\setlength{\emergencystretch}{3em} % prevent overfull lines

\providecommand{\tightlist}{%
  \setlength{\itemsep}{0pt}\setlength{\parskip}{0pt}}



 

\usepackage[]{csquotes}

\IfFileExists{plex-otf.sty}{
  %% Full TeXlive
  % \usepackage[%
  %   % math,
  %   RM={Scale=0.94},SS={Scale=0.94},SScon={Scale=0.94},TT={Scale=MatchLowercase,FakeStretch=0.9},DefaultFeatures={Ligatures=Common}
  % ]{plex-otf}
}{
  %% TinyTeX
  \usepackage{libertine}
}

%%% Load theme
% https://deic.uab.cat/~iblanes/beamer_gallery/
\IfFileExists{beamerthemegotham.sty}{
  %% Full TeXlive
  \usetheme{gotham}
  \gothamset{
    numbering=totalpagenumber,
    parttocframe default=off,
    sectiontocframe default=off,
    subsectiontocframe default=off,
  }
}{
  %% TinyTeX
  \usetheme{Madrid}
}
\makeatletter
\@ifpackageloaded{caption}{}{\usepackage{caption}}
\AtBeginDocument{%
\ifdefined\contentsname
  \renewcommand*\contentsname{Содержание}
\else
  \newcommand\contentsname{Содержание}
\fi
\ifdefined\listfigurename
  \renewcommand*\listfigurename{Список иллюстраций}
\else
  \newcommand\listfigurename{Список иллюстраций}
\fi
\ifdefined\listtablename
  \renewcommand*\listtablename{Список таблиц}
\else
  \newcommand\listtablename{Список таблиц}
\fi
\ifdefined\figurename
  \renewcommand*\figurename{Рисунок}
\else
  \newcommand\figurename{Рисунок}
\fi
\ifdefined\tablename
  \renewcommand*\tablename{Таблица}
\else
  \newcommand\tablename{Таблица}
\fi
}
\@ifpackageloaded{float}{}{\usepackage{float}}
\floatstyle{ruled}
\@ifundefined{c@chapter}{\newfloat{codelisting}{h}{lop}}{\newfloat{codelisting}{h}{lop}[chapter]}
\floatname{codelisting}{Список}
\newcommand*\listoflistings{\listof{codelisting}{Листинги}}
\makeatother
\makeatletter
\makeatother
\makeatletter
\@ifpackageloaded{caption}{}{\usepackage{caption}}
\@ifpackageloaded{subcaption}{}{\usepackage{subcaption}}
\makeatother

\usepackage{bookmark}
\IfFileExists{xurl.sty}{\usepackage{xurl}}{} % add URL line breaks if available
\urlstyle{same}
\hypersetup{
  pdftitle={Первый этап реализации индивидуального проекта},
  pdfauthor={Карпухин Клим},
  pdflang={ru-RU},
  hidelinks,
  pdfcreator={LaTeX via pandoc}}


\title{Первый этап реализации индивидуального проекта}
\subtitle{Размещение на GitHub Pages заготовки для персонального сайта}
\author{Карпухин Клим}
\date{2026-03-07}

\begin{document}
\frame{\titlepage}


\section{Содержание}\label{ux441ux43eux434ux435ux440ux436ux430ux43dux438ux435}

\begin{frame}{Содержание}
\begin{enumerate}[<+->]
\tightlist
\item
  Информация о докладчике
\item
  Цель работы
\item
  Задание
\item
  Теоретическое введение
\item
  Ход работы (этапы, скриншоты)
\item
  Выводы
\end{enumerate}
\end{frame}

\section{Информация}\label{ux438ux43dux444ux43eux440ux43cux430ux446ux438ux44f}

\begin{frame}{Докладчик}
\phantomsection\label{ux434ux43eux43aux43bux430ux434ux447ux438ux43a}
\begin{columns}[c]
\begin{column}{0.65\linewidth}
\begin{itemize}[<+->]
\tightlist
\item
  \textbf{Карпухин Клим}
\item
  Российский университет дружбы народов
\item
  Email:
  \href{mailto:1032255580@rudn.ru}{\nolinkurl{1032255580@rudn.ru}}
\item
  Роли: студент (индивидуальный проект)
\end{itemize}
\end{column}

\begin{column}{0.35\linewidth}

\includegraphics[width=0.9\linewidth,height=\textheight,keepaspectratio]{./image/dokladchik.png}
\end{column}
\end{columns}
\end{frame}

\section{Вводная
часть}\label{ux432ux432ux43eux434ux43dux430ux44f-ux447ux430ux441ux442ux44c}

\begin{frame}{Цель работы}
\phantomsection\label{ux446ux435ux43bux44c-ux440ux430ux431ux43eux442ux44b}
\begin{itemize}[<+->]
\tightlist
\item
  Создать базовую версию персонального веб-сайта с помощью генератора
  статических сайтов Hugo.
\item
  Развернуть сайт на платформе GitHub Pages.
\item
  Настроить автоматическое обновление сайта при изменениях в репозитории
  через GitHub Actions.
\end{itemize}
\end{frame}

\begin{frame}[fragile]{Задание}
\phantomsection\label{ux437ux430ux434ux430ux43dux438ux435}
\begin{enumerate}[<+->]
\tightlist
\item
  Подготовить рабочее окружение: установить необходимое программное
  обеспечение.
\item
  Загрузить шаблон темы для будущего сайта.
\item
  Создать удалённый репозиторий на GitHub и связать его с локальной
  версией проекта.
\item
  Настроить параметр \texttt{baseURL} для корректной работы сайта на
  GitHub Pages.
\item
  Опубликовать заготовку сайта, используя механизмы GitHub Pages и
  GitHub Actions.
\end{enumerate}
\end{frame}

\begin{frame}[fragile]{Теоретическое введение}
\phantomsection\label{ux442ux435ux43eux440ux435ux442ux438ux447ux435ux441ux43aux43eux435-ux432ux432ux435ux434ux435ux43dux438ux435}
\begin{itemize}[<+->]
\tightlist
\item
  \textbf{GitHub Pages} --- сервис для размещения статических
  веб-страниц из репозиториев GitHub.
\item
  \textbf{Hugo} --- генератор статических сайтов, превращающий исходные
  файлы в готовый HTML, CSS и JavaScript.
\item
  \textbf{GitHub Actions} --- автоматизация сборки и публикации: при
  обновлении ветки \texttt{main} запускается процесс, который собирает
  сайт и размещает его на Pages.
\item
  \textbf{Wowchemy (Hugo Academic Theme)} --- стартовый шаблон на основе
  Hugo. Для обработки стилей через TailwindCSS необходимы Node.js и
  pnpm.
\end{itemize}
\end{frame}

\section{Ход работы --- подготовка
окружения}\label{ux445ux43eux434-ux440ux430ux431ux43eux442ux44b-ux43fux43eux434ux433ux43eux442ux43eux432ux43aux430-ux43eux43aux440ux443ux436ux435ux43dux438ux44f}

\begin{frame}{Этап 1: Установка необходимого ПО}
\phantomsection\label{ux44dux442ux430ux43f-1-ux443ux441ux442ux430ux43dux43eux432ux43aux430-ux43dux435ux43eux431ux445ux43eux434ux438ux43cux43eux433ux43e-ux43fux43e}
\begin{itemize}[<+->]
\tightlist
\item
  Установка Git, Go и Node.js через терминал.
\end{itemize}

\begin{figure}

\centering{

\includegraphics[width=0.3\linewidth,height=\textheight,keepaspectratio]{image/1.png}

}

\caption{\label{fig-001}Установка Git, Go и Node.js}

\end{figure}%
\end{frame}

\begin{frame}{Этап 2: Установка Hugo Extended}
\phantomsection\label{ux44dux442ux430ux43f-2-ux443ux441ux442ux430ux43dux43eux432ux43aux430-hugo-extended}
\begin{itemize}[<+->]
\tightlist
\item
  Установка расширенной версии Hugo (Hugo Extended) для поддержки
  SCSS/SASS и TailwindCSS.
\end{itemize}

\begin{figure}

\centering{

\includegraphics[width=0.3\linewidth,height=\textheight,keepaspectratio]{image/2.png}

}

\caption{\label{fig-002}Установка Hugo Extended}

\end{figure}%
\end{frame}

\section{Ход работы --- получение
шаблона}\label{ux445ux43eux434-ux440ux430ux431ux43eux442ux44b-ux43fux43eux43bux443ux447ux435ux43dux438ux435-ux448ux430ux431ux43bux43eux43dux430}

\begin{frame}[fragile]{Этап 3: Клонирование шаблона и установка pnpm}
\phantomsection\label{ux44dux442ux430ux43f-3-ux43aux43bux43eux43dux438ux440ux43eux432ux430ux43dux438ux435-ux448ux430ux431ux43bux43eux43dux430-ux438-ux443ux441ux442ux430ux43dux43eux432ux43aux430-pnpm}
\begin{itemize}[<+->]
\tightlist
\item
  Создание рабочей директории, клонирование репозитория-шаблона
  \texttt{starter-hugo-academic}, установка pnpm.
\end{itemize}

\begin{figure}

\centering{

\includegraphics[width=0.3\linewidth,height=\textheight,keepaspectratio]{image/3.png}

}

\caption{\label{fig-003}Клонирование шаблона и установка pnpm}

\end{figure}%
\end{frame}

\begin{frame}{Этап 4: Запуск локального сервера}
\phantomsection\label{ux44dux442ux430ux43f-4-ux437ux430ux43fux443ux441ux43a-ux43bux43eux43aux430ux43bux44cux43dux43eux433ux43e-ux441ux435ux440ux432ux435ux440ux430}
\begin{itemize}[<+->]
\tightlist
\item
  Запуск локального сервера разработки Hugo для проверки внешнего вида.
\end{itemize}

\begin{figure}

\centering{

\includegraphics[width=0.3\linewidth,height=\textheight,keepaspectratio]{image/4.png}

}

\caption{\label{fig-004}Запуск локального сервера Hugo}

\end{figure}%
\end{frame}

\begin{frame}[fragile]{Этап 5: Просмотр локальной версии}
\phantomsection\label{ux44dux442ux430ux43f-5-ux43fux440ux43eux441ux43cux43eux442ux440-ux43bux43eux43aux430ux43bux44cux43dux43eux439-ux432ux435ux440ux441ux438ux438}
\begin{itemize}[<+->]
\tightlist
\item
  Сайт открылся в браузере по адресу \texttt{http://localhost:1313/} ---
  базовая структура отображается корректно.
\end{itemize}

\begin{figure}

\centering{

\includegraphics[width=0.3\linewidth,height=\textheight,keepaspectratio]{image/5.png}

}

\caption{\label{fig-005}Просмотр локальной версии сайта}

\end{figure}%
\end{frame}

\section{Ход работы --- создание
репозитория}\label{ux445ux43eux434-ux440ux430ux431ux43eux442ux44b-ux441ux43eux437ux434ux430ux43dux438ux435-ux440ux435ux43fux43eux437ux438ux442ux43eux440ux438ux44f}

\begin{frame}[fragile]{Этап 6: Создание репозитория на GitHub}
\phantomsection\label{ux44dux442ux430ux43f-6-ux441ux43eux437ux434ux430ux43dux438ux435-ux440ux435ux43fux43eux437ux438ux442ux43eux440ux438ux44f-ux43dux430-github}
\begin{itemize}[<+->]
\tightlist
\item
  Создание нового публичного репозитория \texttt{my-academic-site} на
  GitHub для хранения исходного кода.
\end{itemize}

\begin{figure}

\centering{

\includegraphics[width=0.3\linewidth,height=\textheight,keepaspectratio]{image/6.png}

}

\caption{\label{fig-006}Создание нового репозитория на GitHub}

\end{figure}%
\end{frame}

\begin{frame}[fragile]{Этап 7: Настройка параметра \texttt{baseURL}}
\phantomsection\label{ux44dux442ux430ux43f-7-ux43dux430ux441ux442ux440ux43eux439ux43aux430-ux43fux430ux440ux430ux43cux435ux442ux440ux430-baseurl}
\begin{itemize}[<+->]
\tightlist
\item
  В файле \texttt{hugo.yaml} указан
  \texttt{baseURL\ =\ https://kreazot316.github.io/my-academic-site/}.
\end{itemize}

\begin{figure}

\centering{

\includegraphics[width=0.3\linewidth,height=\textheight,keepaspectratio]{image/7.png}

}

\caption{\label{fig-007}Настройка baseURL в конфигурации}

\end{figure}%
\end{frame}

\begin{frame}[fragile]{Этап 8: Подготовка контента для стабильной
сборки}
\phantomsection\label{ux44dux442ux430ux43f-8-ux43fux43eux434ux433ux43eux442ux43eux432ux43aux430-ux43aux43eux43dux442ux435ux43dux442ux430-ux434ux43bux44f-ux441ux442ux430ux431ux438ux43bux44cux43dux43eux439-ux441ux431ux43eux440ux43aux438}
\begin{itemize}[<+->]
\tightlist
\item
  Временное отключение раздела блога (перемещение \texttt{content/blog}
  в \texttt{content-disabled/blog}) для исключения ошибок.
\end{itemize}

\begin{figure}

\centering{

\includegraphics[width=0.3\linewidth,height=\textheight,keepaspectratio]{image/8.png}

}

\caption{\label{fig-008}Отключение демо-блога}

\end{figure}%
\end{frame}

\section{Ход работы --- публикация
проекта}\label{ux445ux43eux434-ux440ux430ux431ux43eux442ux44b-ux43fux443ux431ux43bux438ux43aux430ux446ux438ux44f-ux43fux440ux43eux435ux43aux442ux430}

\begin{frame}{Этап 9: Отправка проекта на GitHub}
\phantomsection\label{ux44dux442ux430ux43f-9-ux43eux442ux43fux440ux430ux432ux43aux430-ux43fux440ux43eux435ux43aux442ux430-ux43dux430-github}
\begin{itemize}[<+->]
\tightlist
\item
  Инициализация Git-репозитория, добавление файлов, коммит и отправка в
  удалённый репозиторий.
\end{itemize}

\begin{figure}

\centering{

\includegraphics[width=0.3\linewidth,height=\textheight,keepaspectratio]{image/9.png}

}

\caption{\label{fig-009}Отправка локального проекта на GitHub}

\end{figure}%
\end{frame}

\begin{frame}{Этап 10: Настройка GitHub Pages и GitHub Actions}
\phantomsection\label{ux44dux442ux430ux43f-10-ux43dux430ux441ux442ux440ux43eux439ux43aux430-github-pages-ux438-github-actions}
\begin{itemize}[<+->]
\tightlist
\item
  В настройках репозитория выбран источник публикации \textbf{GitHub
  Actions}.
\end{itemize}

\begin{figure}

\centering{

\includegraphics[width=0.3\linewidth,height=\textheight,keepaspectratio]{image/10.png}

}

\caption{\label{fig-010}Выбор источника GitHub Actions}

\end{figure}%
\end{frame}

\begin{frame}[fragile]{Этап 11: Создание workflow-файла}
\phantomsection\label{ux44dux442ux430ux43f-11-ux441ux43eux437ux434ux430ux43dux438ux435-workflow-ux444ux430ux439ux43bux430}
\begin{itemize}[<+->]
\tightlist
\item
  Создан файл \texttt{.github/workflows/hugo.yml} с описанием шагов
  сборки и деплоя.
\end{itemize}

\begin{figure}

\centering{

\includegraphics[width=0.3\linewidth,height=\textheight,keepaspectratio]{image/11.png}

}

\caption{\label{fig-011}Workflow для автоматической сборки и деплоя}

\end{figure}%
\end{frame}

\begin{frame}{Этап 12: Отправка workflow в репозиторий}
\phantomsection\label{ux44dux442ux430ux43f-12-ux43eux442ux43fux440ux430ux432ux43aux430-workflow-ux432-ux440ux435ux43fux43eux437ux438ux442ux43eux440ux438ux439}
\begin{itemize}[<+->]
\tightlist
\item
  Добавление, коммит и отправка workflow-файла на GitHub.
\end{itemize}

\begin{figure}

\centering{

\includegraphics[width=0.3\linewidth,height=\textheight,keepaspectratio]{image/12.png}

}

\caption{\label{fig-012}Отправка workflow в репозиторий}

\end{figure}%
\end{frame}

\begin{frame}[fragile]{Этап 13: Проверка результата}
\phantomsection\label{ux44dux442ux430ux43f-13-ux43fux440ux43eux432ux435ux440ux43aux430-ux440ux435ux437ux443ux43bux44cux442ux430ux442ux430}
\begin{itemize}[<+->]
\tightlist
\item
  Сайт доступен по адресу
  \texttt{https://kreazot316.github.io/my-academic-site/} и выглядит так
  же, как локальная версия.
\end{itemize}

\begin{figure}

\centering{

\includegraphics[width=0.3\linewidth,height=\textheight,keepaspectratio]{image/13.png}

}

\caption{\label{fig-013}Финальный вид опубликованного сайта}

\end{figure}%
\end{frame}

\section{Выводы}\label{ux432ux44bux432ux43eux434ux44b}

\begin{frame}{Выводы}
\phantomsection\label{ux432ux44bux432ux43eux434ux44b-1}
\begin{itemize}[<+->]
\tightlist
\item
  Выполнена начальная настройка проекта: установлено ПО, загружен шаблон
  Hugo, создан удалённый репозиторий.
\item
  Настроена автоматическая сборка и публикация с помощью GitHub Actions.
\item
  Заготовка персонального сайта доступна в интернете и готова к
  дальнейшему наполнению и модификации.
\end{itemize}
\end{frame}




\end{document}
